\thispagestyle{empty}
\vspace*{25mm}
\noindent{\sffamily\bfseries\Large Abstract\par}
\vspace{6mm}
\noindent
Amphibious quadruped robots can swim by moving their legs, but a leg designed for walking makes a poor
propulsor. This thesis asks which propulsion principle suits the legs of the quadruped BODY2 and how each
propulsor should move. Drag-based paddling with a folding fan paddle is compared with lift-based flapping of a
small fluke fitted to the same leg. Three simulation tools with different roles are used. First, a reduced-order
MuJoCo model with Morison-type link forces and an actuator speed limit is combined with CMA-ES gait
optimization. Whenever the model contains lift and lift is not suppressed, every feasible optimized gait is lift-propelled, and at equal speed the lift gaits
need one half to one third of the power of the best gaits of a drag-only model. Second, the single leg is simulated
with overset-grid CFD in OpenFOAM. For the fan paddle, a trapezoidal velocity law, a 12-mm folded width and
closing the fan at the start of the deceleration raise the thrust from \SI{23.4}{\milli\newton} for the robot's
current gait to \SIrange{73}{83}{\milli\newton} per leg at rest; the range bounds the change of added mass at the
fold, which the two-state model of the folding fan does not resolve. The paddle loses its thrust by about
\SI{0.30}{\metre\per\second}, because its recovery drag and the cost of opening the fan in the oncoming flow grow
while its power-stroke thrust decays. The fluke produces the same thrust per area at rest and loses it more slowly with speed.
Thrust and efficiency cross over at about \SI{0.24}{} and \SI{0.22}{\metre\per\second}. Scheduling the fluke pitch with speed keeps the angle of attack near
\SI{20}{\degree} and raises the Froude efficiency at \SI{0.30}{\metre\per\second} from 0.24 to 0.40. A
quasi-steady self-propulsion balance predicts a cruising speed of \SI{0.30}{\metre\per\second} for the fluke and
\SIrange{0.26}{0.27}{\metre\per\second} for the optimized paddle, with 43\,\% less power for the fluke with
speed-scheduled pitch. A first hardware test with flukes at \SI{2}{\hertz} is 1.3 to 1.6 times slower than predicted, so these
speeds are upper bounds for the present robot. Third, the real-time solver Flare Live supplies the reference fluke motion, but its lattice does not resolve the
foil and it underpredicts the fluke thrust, so it is used for trends only. The CFD set-up
reproduces the thrust coefficient of a flapping-foil experiment to within 8\,\%.

\vspace{12mm}
\noindent{\sffamily\bfseries\Large Kurzfassung\par}
\vspace{6mm}
\noindent\begingroup\sisetup{locale=DE}
Amphibische Vierbeiner können schwimmen, indem sie ihre Beine bewegen, doch ein zum Laufen ausgelegtes Bein
ist ein schlechter Antrieb. Diese Arbeit untersucht, welches Antriebsprinzip zu den Beinen des Vierbeiners
BODY2 passt und wie sich der jeweilige Antrieb bewegen sollte. Verglichen werden widerstandsbasiertes Paddeln
mit einem faltbaren Fächerpaddel und auftriebsbasiertes Schlagen mit einer kleinen Flosse am selben Bein. Dazu
werden drei Simulationswerkzeuge mit unterschiedlichen Aufgaben eingesetzt. Zunächst wird ein Modell reduzierter Ordnung in
MuJoCo mit Morison-Kräften je Glied und begrenzter Stellgeschwindigkeit der Servos mit einer CMA-ES-Gangoptimierung kombiniert.
Sobald das Modell Auftrieb enthält und dieser nicht unterdrückt wird, nutzen alle zulässigen optimierten Gangarten Auftrieb, und bei gleicher Geschwindigkeit benötigen
sie nur die Hälfte bis ein Drittel der Leistung der besten Gangarten eines reinen Widerstandsmodells. Anschließend
wird das Einzelbein mit Overset-CFD in OpenFOAM simuliert. Für das Fächerpaddel erhöhen ein trapezförmiges
Geschwindigkeitsprofil, eine gefaltete Breite von 12\,mm und das Schließen des Fächers zu Beginn der
Verzögerung den Schub im Stand von \SI{23.4}{\milli\newton} (aktueller Gang des Roboters) auf
\SIrange{73}{83}{\milli\newton} pro Bein; die Spanne erfasst die Änderung der hydrodynamischen Masse beim Falten, die das
Zwei-Zustands-Modell des Fächers nicht auflöst. Bei etwa \SI{0.30}{\metre\per\second} verliert das Paddel seinen Schub, weil
Rückholwiderstand und der Aufwand zum Öffnen des Fächers in der Anströmung wachsen, während der Schub im Arbeitshub abnimmt.
Die Flosse erzeugt im Stand denselben Schub pro Fläche und verliert ihn mit der Geschwindigkeit langsamer. Schub und Wirkungsgrad kreuzen sich bei etwa
\SI{0.24}{} bzw. \SI{0.22}{\metre\per\second}.
Eine geschwindigkeitsabhängige Anstellung hält den Anstellwinkel bei etwa \SI{20}{\degree} und steigert den
Froude-Wirkungsgrad bei \SI{0.30}{\metre\per\second} von 0,24 auf 0,40. Eine quasistationäre
Eigenantriebsbilanz sagt für die Flosse \SI{0.30}{\metre\per\second} und für das optimierte Paddel
\SIrange{0.26}{0.27}{\metre\per\second} voraus, wobei die Flosse mit geschwindigkeitsabhängiger Anstellung 43\,\% weniger
Leistung benötigt. Ein erster Schwimmversuch des Roboters mit Flossen ist bei \SI{2}{\hertz} 1,3- bis 1,6-mal langsamer als
vorhergesagt; die vorhergesagten Geschwindigkeiten sind daher Obergrenzen für den aktuellen Roboter. Der Echtzeitlöser Flare Live
liefert die Referenzbewegung der Flosse, löst das Profil jedoch nicht auf und unterschätzt den Flossenschub; er dient daher nur zur
Einordnung von Trends. Der CFD-Aufbau gibt den Schubbeiwert eines Schlagflügelexperiments auf 8\,\% genau wieder.\endgroup
\cleardoublepage
